\documentclass[11pt,a4paper]{article}

\usepackage[margin=2.6cm]{geometry}
\usepackage{amsmath,amssymb,amsthm}
\usepackage{hyperref}
\usepackage{xurl}
\usepackage{booktabs}
\usepackage{enumitem}
\usepackage{natbib}
\bibliographystyle{plainnat}

\setlength{\emergencystretch}{3em}

\newtheorem{theorem}{Theorem}[section]
\newtheorem{proposition}[theorem]{Proposition}
\newtheorem{lemma}[theorem]{Lemma}
\newtheorem{corollary}[theorem]{Corollary}
\newtheorem{axiom}[theorem]{Axiom}
\theoremstyle{remark}
\newtheorem{remark}[theorem]{Remark}

\newcommand{\Q}{\mathbb{Q}}
\newcommand{\Z}{\mathbb{Z}}
\newcommand{\N}{\mathbb{N}}
\newcommand{\fract}[1]{\left\{#1\right\}}
\newcommand{\lean}[1]{\texttt{#1}}
\newcommand{\gap}{\operatorname{gap}}

\title{A Finiteness Theorem for the Algebraic--Pair Branch of the Base $\tfrac32$ Problem}
\author{Su MingKai}
\date{2026}

\begin{document}

\begin{abstract}
We give a complete proof of the algebraic--pair finiteness theorem for the base $\tfrac32$ Mahler
problem, following the program of Stephan.  For an algebraic irrational multiplier $\delta$ there
are only finitely many ``violating'' pairs $(a,c)$, $2\le a<c$, with distinct gaps $c-a$, for
which $\delta\bigl((3/2)^{c}-(3/2)^{a}\bigr)$ is within $\vartheta^{c}$ of an integer.  The only
external input is Schlickewei's $p$-adic subspace theorem on $\Q^{3}$, recorded in the formalisation
as a single declared axiom; everything else is proved.  The argument splits the rational linear
relation produced by the subspace theorem into two complementary cases: $C=0$, excluded by strict
monotonicity of $s\mapsto(3/2)^{s}$ together with injectivity of the gap map; and $C\neq0$, which
forces the algebraic irrational $\delta$ to be rational.  We formalise the whole reduction in Lean~4
with Mathlib.  We do \emph{not} claim Mahler's original conjecture; this is the algebraic--pair
branch of a conditional strategy and is labelled as such.
\end{abstract}

\tableofcontents

\section{Introduction}

Mahler's problem asks for the distribution of $\fract{\xi(3/2)^{n}}$ and is open; the state of the
art is surveyed in \citep{stephan2026subspace,stephan2026minimal,stephan2026third}.  A point
$\xi$ is a \emph{Z-number} when $\fract{\xi(3/2)^{n}}\in[0,\tfrac12)$ for every $n$; Strauch
\citep{strauch1997} constructed Z-numbers with prescribed maps, and the combinatorial side is
developed in \citep{lagarias2003,andrieu2026,cassaigne2026}.  Recent work of Cook
\citep{cook2026} pushes the arithmetic envelope of base $\tfrac32$ congruences, while the
denominator and equidistribution barriers are those of \citep{calegari2025,elmv}.

Stephan's program splits the surviving ``minimal words'' into a rational case and an
\emph{algebraic--pair} case \citep{stephan2026subspace,stephan2026third}.  The purpose of this
paper is to give the algebraic--pair branch in full: a self-contained proof that, for an algebraic
irrational multiplier, only finitely many pairs can carry a near-integer value.  This is the content
of Theorem~\ref{thm:main}; Sections~\ref{sec:setup}--\ref{sec:main} develop it, and
Section~\ref{sec:algpair-C0}--\ref{sec:algpair-Cne0} treat the two cases that the proof separates.

\paragraph{Scope and honesty.}
The only non-elementary input is Schlickewei's $p$-adic subspace theorem on $\Q^{3}$
(Bombieri--Gubler Ch.~7; Evertse's lecture notes), declared in the formalisation as the single axiom
\lean{schlickewei\_subspace\_Q3}.  Every other result in this paper is unconditional and
machine-checked with zero axioms beyond \lean{propext}, \lean{Quot.sound}, \lean{Classical.choice}.
No statement below asserts the Mahler conjecture.

\section{Setup, hypotheses, and the subspace input}
\label{sec:setup}

A \emph{pair} is an element $p=(a,c)\in\N^{2}$ with $a<c$; its \emph{gap} is $\gap(p)=c-a$.  We
work with an infinite set $T\subseteq\N^{2}$ of pairs on which the gap map is injective, each pair
satisfies $2\le a<c$, and a fixed integer-valued function $n\colon\N^{2}\to\Z$ approximates a
multiple of $\delta$ to exponential accuracy.

\begin{definition}[Approximation hypothesis]
\label{def:hyp}
Let $\delta\in\R$, $\vartheta\in\Q$ with $0<\vartheta<1$, and $n\colon\N^{2}\to\Z$.  The pair set
$T$ satisfies $(\mathrm{H})$ if
\[
  \bigl|\delta\bigl((3/2)^{c}-(3/2)^{a}\bigr)-n(p)\bigr|\le\vartheta^{c}
  \qquad\text{for all }p=(a,c)\in T .
\]
\end{definition}

The engine of the proof is the following classical theorem, used verbatim.

\begin{axiom}[Schlickewei $p$-adic subspace theorem on $\Q^{3}$]
\label{ax:schlickewei}
Let $\delta\in\R$ be algebraic over $\Q$, let $0<\vartheta<1$ be rational, and let
$T\subseteq\N^{2}$ be infinite with $2\le a<c$ for all $(a,c)\in T$ and $(\mathrm{H})$ satisfied by
some $n$.  Then there is an infinite subset $T'\subseteq T$ and a nonzero rational triple
$(A,B,C)\neq(0,0,0)$ with
\[
  A\,(3/2)^{c}+B\,(3/2)^{a}+C\,n(p)=0
  \qquad\text{for all }p=(a,c)\in T'.
\]
\end{axiom}

The three Archimedean, $2$-adic and $3$-adic product estimates of the task card cancel exactly to
produce the height bound $\vartheta^{c}$ that the subspace theorem requires; we take the conclusion
as the single black box.

\section{The case $C=0$}
\label{sec:algpair-C0}

We first exclude the degenerate relation, in which only $(3/2)^{a}$ and $(3/2)^{c}$ appear.

\begin{lemma}[Strict monotonicity]
\label{lem:mono}
The map $s\mapsto(3/2)^{s}$ is strictly increasing on $\N$.
\end{lemma}

\begin{proof}
Since $(3/2)>1$, $m<n$ implies $(3/2)^{m}<(3/2)^{n}$.
Machine-checked inside the module as \lean{pow\_three\_halves\_strictMono}.
\end{proof}

\begin{lemma}[An infinite set has two distinct points]
\label{lem:twopoints}
If $S$ is infinite, there exist $p,q\in S$ with $p\neq q$.
\end{lemma}

\begin{proof}
Otherwise $S$ would be a subsingleton, hence finite, contradicting infinitude.
\end{proof}

\begin{theorem}[No $C=0$ relation]
\label{thm:C0}
Let $T'$ be infinite, with injective gap map and $a<c$ on $T'$.  If $A,B\in\Q$, not both zero,
satisfy
\[
  A\,(3/2)^{c}+B\,(3/2)^{a}=0\qquad\text{for all }p=(a,c)\in T',
\]
then contradiction.
\end{theorem}

\begin{proof}
$A\neq0$: if $A=0$ then $B\neq0$ and $B(3/2)^{a}=0$ at any $p\in T'$, impossible since
$(3/2)^{a}>0$.  Dividing the relation by $A(3/2)^{a}$ gives
\[
  (3/2)^{c-a}=-\frac{B}{A}\qquad\text{for all }p=(a,c)\in T'.
\]
Thus $(3/2)^{\gap(p)}$ is the fixed constant $-B/A$ on $T'$.  Taking two distinct points
$p\neq q$ of $T'$ (Lemma~\ref{lem:twopoints}) gives
$(3/2)^{\gap(p)}=(3/2)^{\gap(q)}$, so by Lemma~\ref{lem:mono} the exponents agree,
$\gap(p)=\gap(q)$.  Since the gap map is injective on $T'$, this forces $p=q$, a contradiction.
Machine-checked: \lean{Mahler32Conditional.lemma1\_no\_C\_zero}.
\end{proof}

\section{The case $C\neq0$}
\label{sec:algpair-Cne0}

Now assume the relation genuinely involves $n(p)$.  Solving it for $n(p)$ turns the approximation
hypothesis into an approximation of $\delta$ by a fixed rational, which we then push to the limit.

\begin{lemma}[One-step bound]
\label{lem:single}
Let $a<c$ and suppose
\[
  \bigl|\delta\bigl((3/2)^{c}-(3/2)^{a}\bigr)-\bigl(\alpha(3/2)^{c}+\beta(3/2)^{a}\bigr)\bigr|
  \le\vartheta^{c}
\]
with $\alpha,\beta\in\Q$.  Then, writing $X=3/2$,
\[
  \bigl|\delta-\alpha\bigr|\le\bigl|\delta+\beta\bigr|\,(2/3)^{c-a}+(2\vartheta/3)^{c}.
\]
\end{lemma}

\begin{proof}
Rearranging the hypothesis,
\[
  \bigl|X^{c}(\delta-\alpha)-X^{a}(\delta+\beta)\bigr|\le\vartheta^{c}
\]
because $\delta(X^{c}-X^{a})-(\alpha X^{c}+\beta X^{a})=X^{c}(\delta-\alpha)-X^{a}(\delta+\beta)$.
The reverse triangle inequality gives
$X^{c}|\delta-\alpha|\le X^{a}|\delta+\beta|+\vartheta^{c}$.  Multiplying by $(2/3)^{c}$ and using
$(2/3)^{c}X^{c}=1$,
\[
  |\delta-\alpha|\le|\delta+\beta|\,(2/3)^{c}X^{a}+(2/3)^{c}\vartheta^{c}.
\]
Finally $(2/3)^{c}X^{a}=(2/3)^{c-a}(2/3)^{a}X^{a}=(2/3)^{c-a}$ and
$(2/3)^{c}\vartheta^{c}=(2\vartheta/3)^{c}$.  This is the stated bound.
\end{proof}

\begin{lemma}[The gap is unbounded]
\label{lem:gapunbounded}
If $T'$ is infinite and the gap map is injective on $T'$, then for every $N$ there is $p\in T'$ with
$\gap(p)\ge N$.
\end{lemma}

\begin{proof}
If all gaps were $<N$, then $T'$ would be the union over the finitely many values
$s\in\{0,\dots,N-1\}$ of the set $\{p\in T':\gap(p)=s\}$, each of which has at most one element by
injectivity.  Hence $T'$ would be finite, a contradiction.
Machine-checked inside the module as \lean{gap\_unbounded}.
\end{proof}

\begin{theorem}[No $C\neq0$ relation]
\label{thm:Cne0}
Let $\delta$ be irrational, $0<\vartheta<1$, and $T'$ infinite with injective gap map and $a<c$ on
$T'$.  If $\alpha,\beta\in\Q$ satisfy
\[
  \bigl|\delta\bigl((3/2)^{c}-(3/2)^{a}\bigr)-\bigl(\alpha(3/2)^{c}+\beta(3/2)^{a}\bigr)\bigr|
  \le\vartheta^{c}\qquad\text{for all }p=(a,c)\in T',
\]
then contradiction.
\end{theorem}

\begin{proof}
By Lemma~\ref{lem:single}, for every $p=(a,c)\in T'$,
\[
  |\delta-\alpha|\le|\delta+\beta|\,(2/3)^{\gap(p)}+(2\vartheta/3)^{c}.
\]
Let $\varepsilon>0$.  Since $0<2/3<1$ and $0<2\vartheta/3<2/3<1$, choose $N$ with
$|\delta+\beta|(2/3)^{N}<\varepsilon/2$ and $(2\vartheta/3)^{N}<\varepsilon/2$.  By
Lemma~\ref{lem:gapunbounded} pick $p\in T'$ with $\gap(p)\ge N$; then $c\ge\gap(p)\ge N$, so both
terms are below $\varepsilon/2$ and $|\delta-\alpha|<\varepsilon$.  As $\varepsilon$ was arbitrary,
$|\delta-\alpha|=0$, i.e.\ $\delta=\alpha\in\Q$, contradicting the irrationality of $\delta$.
Machine-checked: \lean{Mahler32Conditional.lemma2\_no\_C\_ne\_zero}.
\end{proof}

\section{The finiteness theorem}
\label{sec:main}

We can now assemble the two cases.

\begin{theorem}[Algebraic--pair finiteness, conditional on
Schlickewei]\label{thm:main}
Let $\delta$ be algebraic over $\Q$ and irrational, let $0<\vartheta<1$ be rational, and let
$T\subseteq\N^{2}$ satisfy $2\le a<c$ on $T$, have injective gap map, and satisfy $(\mathrm{H})$ for
some $n\colon\N^{2}\to\Z$.  Then $T$ is finite.
\end{theorem}

\begin{proof}
Suppose, toward a contradiction, that $T$ is infinite.  Axiom~\ref{ax:schlickewei} yields an
infinite $T'\subseteq T$ and a nonzero rational triple $(A,B,C)$ with
\[
  A\,(3/2)^{c}+B\,(3/2)^{a}+C\,n(p)=0\qquad\text{for all }p=(a,c)\in T'.
\]
\emph{Case $C=0$.}  Then $A(3/2)^{c}+B(3/2)^{a}=0$ on $T'$ with $(A,B)\neq(0,0)$; the hypotheses of
Theorem~\ref{thm:C0} hold ($T'$ infinite, injective gaps, $a<c$), giving a contradiction.

\emph{Case $C\neq0$.}  Set $\alpha=-A/C$ and $\beta=-B/C$, so the relation reads
$n(p)=\alpha(3/2)^{c}+\beta(3/2)^{a}$ on $T'$.  Substituting into $(\mathrm{H})$,
\[
  \bigl|\delta\bigl((3/2)^{c}-(3/2)^{a}\bigr)-\bigl(\alpha(3/2)^{c}+\beta(3/2)^{a}\bigr)\bigr|
  \le\vartheta^{c}\qquad\text{for all }p\in T',
\]
which is exactly the hypothesis of Theorem~\ref{thm:Cne0}; with $\delta$ irrational this is again a
contradiction.

Both cases are impossible, so $T$ is finite.
Machine-checked: \lean{Mahler32Conditional.algebraic\_pair\_branch\_finite}, depending on the single
axiom \lean{schlickewei\_subspace\_Q3}.
\end{proof}

\begin{remark}[Read this before citing]
Theorem~\ref{thm:main} is \emph{conditional} on Axiom~\ref{ax:schlickewei}.  Every unconditional
statement in this paper (Sections~\ref{sec:algpair-C0}--\ref{sec:algpair-Cne0}, and the assembly of
Theorem~\ref{thm:main} apart from the one axiom application) is zero-axiom in the kernel.  The CI
axiom audit lists \lean{schlickewei\_subspace\_Q3} explicitly and excludes it from the unconditional
count.  We make no claim that the Mahler problem is settled.
\end{remark}

\section{What is and is not claimed}
\label{sec:claims}

\begin{itemize}[leftmargin=1.4em]
  \item \textbf{Unconditional and machine-checked.} The $C=0$ exclusion
        (Theorem~\ref{thm:C0}), the one-step bound (Lemma~\ref{lem:single}), and the $C\neq0$
        exclusion (Theorem~\ref{thm:Cne0}).
  \item \textbf{Conditional and machine-checked.} Theorem~\ref{thm:main}, modulo the single
        declared axiom \lean{schlickewei\_subspace\_Q3}.
  \item \textbf{Human-level, not formalised.} The $p$-adic product estimates that produce the
        height bound $\vartheta^{c}$ and hence the hypothesis of the subspace theorem.  These are
        taken from the literature \citep{stephan2026subspace} and are not re-proved here.
  \item \textbf{Not claimed at all.} Mahler's original conjecture, the unbounded-denominators
        conjecture, and any unconditional transcendence of a Z-number.
\end{itemize}

\section{Formalisation traceability}

All Lean names below live in the repository and are exercised by \texttt{scripts/verify\_all.sh}.
The axiom audit prints only \lean{propext}, \lean{Quot.sound}, \lean{Classical.choice} for the
unconditional results; the conditional theorem additionally reports \lean{schlickewei\_subspace\_Q3}.

\begin{table}[h]
\centering
\small
\begin{tabular}{@{}lll@{}}
\toprule
Result & Lean declaration & Status\\
\midrule
$C=0$ exclusion & \lean{Mahler32Conditional.lemma1\_no\_C\_zero} & proved\\
One-step bound & inside \lean{Mahler32Conditional.lemma2\_no\_C\_ne\_zero} & proved\\
$C\neq0$ exclusion & \lean{Mahler32Conditional.lemma2\_no\_C\_ne\_zero} & proved\\
Finiteness (cond.) & \lean{Mahler32Conditional.algebraic\_pair\_branch\_finite} & axiom-dependent\\
Subspace input & \lean{Mahler32Conditional.schlickewei\_subspace\_Q3} & declared axiom\\
\bottomrule
\end{tabular}
\caption{Theorem-to-Lean map for Paper~1.}
\end{table}

\bibliography{refs}

\end{document}
